\documentclass[]{article}

\usepackage{amsmath}
\usepackage{multirow}
\usepackage{natbib}
\usepackage{graphicx} 


\begin{document}

Our high-throughput screening pipeline evaluated 889 unique oxide-carbonate reaction pairs, filtered for thermodynamic stability, to identify materials that optimally balance the competing requirements for an ideal CO₂ sorbent. The analysis systematically quantified the thermodynamic driving force for capture, the mechanical strain from volume expansion, and the thermal resistance to sintering. The results reveal clear structure-property relationships and identify a new class of complex oxide materials that overcome the limitations of simple binary oxides.

\begin{figure*}[t]
    \centering
    \includegraphics[width=\textwidth]{../input_files/plots/step_6_multi_panel_analysis_1_1775790796.png}
    \caption{Multifaceted screening of CO₂ sorbent properties. (Top-left) A Pareto analysis of Gibbs free energy ($\Delta$G) versus volume expansion ($\Delta$V) at 700 K demonstrates that most materials are poor sorbents due to extreme thermodynamics or catastrophic volume expansion (> 20\%), isolating a few promising candidates in the optimal region near $\Delta G = 0$ and low $\Delta V$. (Top-right) Box plots of reaction enthalpy ($\Delta$H) show that alkali and alkaline earth metals (e.g., K, Na) are highly exothermic and non-regenerable, while mixed-metal polyanionic systems (e.g., Li-P-V, Fe-Li-P) exhibit the moderate thermodynamics ideal for reversible cycling. (Bottom-left) The distribution of the Tamman temperature operating margin reveals that many candidates possess high thermal stability and resistance to sintering. (Bottom-right) Isostructural substitution analysis establishes a design rule where cation substitution systematically tunes reaction enthalpy, with less electropositive metals like Ni and Co destabilizing the carbonate and lowering the regeneration energy.}
    \label{fig:pareto_and_distributions}
\end{figure*}

\subsection{Thermodynamic, mechanical, and thermal landscape}
The overall performance landscape of the 889 candidates is summarized in Figure \ref{fig:pareto_and_distributions}. The thermodynamic viability of each reaction was assessed by calculating the Gibbs free energy ($\Delta G$) at 700 K and a CO₂ partial pressure of 0.15 bar. As shown in the top-right panel of Figure \ref{fig:pareto_and_distributions}, the reaction enthalpies ($\Delta H$) cluster strongly by the primary metal cation. Alkali and alkaline earth metals, such as K, Na, and Ba, exhibit highly exothermic reactions (mean $\Delta H$ of -4.22, -4.20, and -3.69 eV/CO₂, respectively), resulting in extremely negative $\Delta G$ values. While these materials have a strong affinity for CO₂, their resulting carbonates are too stable for energy-efficient thermal regeneration. In contrast, many simple transition metal oxides, such as those of iron, have positive $\Delta G$ values at 700 K (mean of +0.101 eV), rendering them non-reactive under typical flue gas conditions. The most promising materials occupy an intermediate, "Goldilocks" thermodynamic window. Mixed-metal and polyanionic systems, such as Li-P-V and Li-Mn-P, exhibit mean $\Delta G$ values of -0.329 eV and -0.475 eV, respectively. This moderation is achieved by incorporating electronegative polyanions (e.g., phosphates), which reduce the basicity of the lattice oxygen and tune the reaction thermodynamics for reversible cycling.

Beyond thermodynamics, long-term durability requires resistance to thermal degradation. We estimated each material's resistance to sintering using the Tamman temperature ($T_T$), which marks the onset of significant lattice diffusion. The distribution of the operating margin ($T_T - 700$ K), shown in the bottom-left panel of Figure \ref{fig:pareto_and_distributions}, highlights several highly robust candidates. The calcium phosphate system (Ca₁₀P₆O₂₅) displays the highest resistance, with an operating margin of 1028.9 K. Other complex oxides, including potassium titanium phosphate (KTiPO₅) and sodium titanium phosphate (NaTiPO₅), also show excellent thermal stability with margins exceeding 930 K. This suggests that rigid polyhedral units (e.g., PO₄, SiO₄, TiO₆) enhance the cohesive energy of the crystal lattice, making these materials resilient to the high temperatures required for cyclic operation.

\subsection{Volumetric expansion as a critical failure mode}
A primary mechanism for the mechanical failure of solid sorbents is the stress induced by volume changes during carbonation and decarbonation cycles. Our analysis reveals that this is a dominant and often disqualifying factor. Figure \ref{fig:volume_expansion} shows the distribution of molar volume expansion ($\Delta V$) for all 889 reaction pairs. The distribution is heavily skewed, with a mean expansion of 74.85\% and some materials expanding by over 200\%. Such large, repeated structural changes lead to particle pulverization, loss of surface area, and a rapid decline in CO₂ capture capacity.

\begin{figure}[h!]
    \centering
    \includegraphics[width=0.5\textwidth]{../input_files/plots/step_3_volume_expansion_histogram_1775790499.png}
    \caption{Distribution of the molar volume expansion ($\Delta$V) upon carbonation for the 889 screened reaction pairs. The histogram reveals a heavily right-skewed distribution with a mean expansion of 74.85\%, indicating that most materials undergo significant structural swelling detrimental to mechanical integrity. The red dashed lines highlight the strict mechanical robustness filter applied ($|\Delta V| \leq 20\%$), visually demonstrating the rarity of dimensionally stable candidates, with only 14 reactions meeting this criterion.}
    \label{fig:volume_expansion}
\end{figure}

To identify mechanically robust candidates, we applied a stringent filter, retaining only materials with an absolute volume change of 20\% or less. As highlighted in Figure \ref{fig:volume_expansion}, this filter drastically reduces the candidate pool to just 14 reaction pairs, underscoring the rarity of dimensionally stable sorbents. The top-left panel of Figure \ref{fig:pareto_and_distributions} further illustrates this, showing that the vast majority of materials lie outside the desired low-$\Delta V$ region. The materials that do satisfy this criterion are predominantly complex, open-framework structures. For instance, the sodium aluminosilicate system (Na₈Al₆Si₆O₂₅) exhibits a near-zero volume change (-0.77\% to 1.17\%), as CO₂ can be accommodated within its existing microporous structure. Similarly, the calcium phosphate system (Ca₁₀P₆O₂₅) shows a minimal expansion of ~1.9\%. These results pinpoint open-framework polyanionic structures as a key motif for designing mechanically durable sorbents.

\subsection{Guiding sorbent design through isostructural substitution}
To move beyond screening and establish rational design principles, we analyzed the effect of cation substitution on reaction thermodynamics within a fixed crystal structure. Using the CaO (Fm-3m) $\rightarrow$ CaCO₃ (R-3c) reaction as a benchmark, we quantified the shift in reaction enthalpy ($\Delta H_{shift}$) upon replacing Ca with other cations, as shown in the bottom-right panel of Figure \ref{fig:pareto_and_distributions}.

The results demonstrate a clear and predictable trend linked to cation electronegativity. Substituting Ca with more electropositive alkaline earth metals like Sr and Ba makes the reaction more exothermic, with $\Delta H_{shift}$ values of -0.402 eV and -0.691 eV, respectively. This increases the stability of the carbonate and raises the regeneration temperature. Conversely, substituting Ca with less electropositive transition metals like Mn, Co, and Ni systematically destabilizes the carbonate, yielding positive $\Delta H_{shift}$ values of +0.719 eV, +0.967 eV, and +0.991 eV. This destabilization significantly lowers the regeneration energy, with Ni substitution reducing the equilibrium temperature by over 600 K relative to the CaO system. This analysis confirms that the M-O bond covalency is a primary tuning parameter; as the bond becomes more covalent, the basicity of the lattice oxygen decreases, weakening its affinity for CO₂. This provides a powerful, predictive rule for fine-tuning the thermodynamics of sorbent materials.

\subsection{Composite ranking and promising polyanionic sorbents}
To identify the most promising candidates overall, we integrated the key performance metrics—thermodynamic reversibility ($\Delta G$), mechanical stability ($\Delta V$), and thermal durability (Tamman margin)—into a single composite score. This holistic ranking moves the focus from materials that excel in one property to those that strike an optimal balance among all three.

The top-ranked candidates are overwhelmingly complex, mixed-metal polyanionic frameworks, representing a significant departure from traditional simple oxides. The highest-scoring material is sodium titanium phosphate (NaTiPO₅), which exhibits a nearly ideal $\Delta G$ of +0.010 eV at 700 K, a moderate volume expansion of 47.8\%, and an exceptional thermal operating margin of 934.2 K. Closely following is potassium titanium phosphate (KTiPO₅), with a $\Delta G$ of -0.003 eV and a margin of 958.3 K. Other promising candidates include lithium vanadium phosphates (e.g., LiVPO₅) and silicates (e.g., Li₂VSiO₅), which also balance near-zero $\Delta G$ values with manageable volume expansions (30-45\%) and high thermal stability.

As visualized in the Pareto plot (Figure \ref{fig:pareto_and_distributions}, top-left), these top candidates occupy the coveted region of near-zero $\Delta G$ and low-to-moderate $\Delta V$. Unlike CaO, which is thermodynamically robust but requires very high regeneration temperatures and suffers from sintering, these polyanionic materials are predicted to cycle reversibly at lower temperatures while resisting both mechanical and thermal degradation. Their discovery provides a data-driven roadmap for the experimental development of a new generation of durable, high-performance solid sorbents for industrial CO₂ capture.

\end{document}
                